\documentclass{article}
\usepackage{amsmath}
\usepackage[utf8]{inputenc}
\usepackage{booktabs}
\usepackage{microtype}
\pagestyle{empty}

\title{closest pair Report}
\author{Author 1 and Author 2}

\begin{document}
  \maketitle

  \section{Results}
  The testing script confirmed that all our solutions were correct. When running the code with the PyPy compiler, it easily handles the largest inputs in a few seconds. The vast majority of the execution time is spent on the initial sorting of the points and the overhead of the recursive function calls. 

  \section{Implementation details}
  I implemented a divide and conquer algorithm. Sorted the points by their x-coordinates. The function then recursively splits the points into left and right halves to find the smallest distance on either side. Finally, we check a narrow strip down the middle to see if any points cross the boundary. 
  
  For data structures I just used standard Python lists of tuples to store the coordinates. Rather than maintaining a second list sorted by y-coordinates throughout the whole program I created a new list for the middle strip and sorted it on the fly using Python's built-in sort (fast c optimized). 
  
  The overall running time is O(n log n).

\end{document} 